\origpage{81}{GRAPH MINORS. XI}
\noindent $r' \subseteq \Delta'$, since $r' \in R(\Gamma')$. Since $\Gamma'$ is confined, $C'$ has length $< 2\theta$; but $U(C') \subseteq \operatorname{ins}(C)$, and so by the first slope axiom, $\operatorname{ins}(C') \subseteq \operatorname{ins}(C)$, that is, $\operatorname{ins}(C') = \Delta'$. It follows that $\Gamma'$ captures $r'$, as required. $\square$

\sectionhead{5. Small Sets}

Again, throughout this section $\Gamma$ is a drawing in a connected surface $\Sigma$ with a slope $\operatorname{ins}$ of order $\theta$. A weighting is a pair $(\Omega, \omega)$, where $\Omega \subseteq V(\Gamma)$, $\omega \ge 1$ is an integer and $|\Omega \cap V(C)| \ge \omega$ for every circuit $C$ of $\Gamma$ of length $\ge 2\theta$. For instance, $(V(\Gamma), 2\theta)$ is a weighting; or if $\Gamma$ is bipartite and $\theta$ is integral, we could take $\Omega$ to be one side of a bipartition and $\omega = \theta$. Throughout this section $(\Omega, \omega)$ is a fixed weighting.

Let $Z \subseteq R(\Gamma)$. We define $N_\Gamma(Z) = N(Z)$ to be the subdrawing of $\Gamma$ formed by those vertices and edges of $\Gamma$ incident both with a member of $Z$ and with a member of $R(\Gamma) - Z$. We say that $Z$ is \textit{small} (with respect to $(\Omega, \omega)$) if $|\Omega \cap V(N(Z))| < \omega$, and $N(Z)$ captures $r$ for each $r \in Z$. This section is devoted to proving certain properties of small sets, which will be central when we come to establish the connection between slopes and tangles.

We begin with the following lemma.

\stmt
(5.1) \textit{Let $Z \subseteq R(\Gamma)$, and let $\Gamma' \subseteq \Gamma$. Let $r'$ be a region of $\Gamma'$. If there are regions $r_1 \in Z$, $r_2 \in R(\Gamma) - Z$ with $r_1, r_2 \subseteq r'$, then there are such regions $r_1, r_2$ which in addition are incident with some edge $e$ of $\Gamma$ with $e \notin E(\Gamma')$.}

\textit{Proof.} Since $r_1, r_2 \subseteq r'$ and $r'$ is a region of $\Gamma'$ there is a sequence
\begin{equation*}
r_1 = a_1, a_2, \dots, a_s = r_2
\end{equation*}
of regions $a_1, a_2, \dots, a_s$ of $\Gamma$ with $a_1, a_2, \dots, a_s \subseteq r'$, such that for $1 \le i < s$ there is an edge $e_i$ of $\Gamma$ incident with both $a_i$ and $a_{i+1}$, not in $E(\Gamma')$. Choose $i$ maximum such that $a_i \in Z$ (this is possible since $a_1 \in Z$); then $i \ne s$ since $a_s \notin Z$, and $a_i$, $a_{i+1}$ satisfy the theorem. $\square$

We apply (5.1) to deduce the following.

\stmt
(5.2) \textit{If $Z \subseteq R(\Gamma)$ and $|\Omega \cap V(N(Z))| < \omega$, then exactly one of $Z$, $R(\Gamma) - Z$ is small.}

\textit{Proof.} Certainly $N(Z)$ is confined, from the definition of a weighting, and so by (4.4) there is a unique region $r_0$ of $N(Z)$ not captured by $N(Z)$. Thus for $r \in R(\Gamma)$, $N(Z)$ captures $r$ if and only if $r \not\subseteq r_0$. It follows that at most one of $Z$, $R(\Gamma) - Z$ is small. Suppose, for a contradiction, that neither is small, so that there exist $r_1 \in Z$, $r_2 \in R(\Gamma) - Z$ with $r_1, r_2 \subseteq r_0$. By (5.1),
